\section{Conclusion}\label{sec:conclusion}
We introduced Conformal Reliability Routing, a decoupled framework for low-storage few-shot industrial anomaly detection. CRR preserves frozen DINOv2 PCA/subspace residuals for anomaly ranking and adds separate calibration and conformal routes for probability-style reliability. The main empirical results are threefold. First, LOIO conformal reliability substantially reduces calibration error on full VisA and full MVTec corruption benchmarks (overall ECE 0.0766 and 0.0684) while leaving the raw ranking unchanged, and it does so against the complete standard toolbox---Vector and Shift-Aware Platt as well as temperature scaling, isotonic regression, and histogram binning fit on the identical label-free calibration set---with paired Wilcoxon support in every comparison except Shift-Aware Platt at VisA $k{=}8$. Second, the conformal promise itself is audited rather than assumed: an attainable-alpha analysis makes the structural silence of $k$-shot conformal alarms below $1/(k{+}1)$ explicit, and exact discrete uniformity tests show that the first attainable target-only operating point is conservative on VisA but anti-conservative on corrupted MVTec. Third, SC3R, a source-validated cross-category thresholding scheme, passes a pre-registered gate on all 15 MVTec classes at $k{=}4$---tracking nominal false-alarm rates almost exactly below and at the floor without target anomaly labels, with CI-supported power gains under every corruption condition---replicates on all 12 VisA classes, and transfers conservatively when the source archive comes from a different dataset entirely, trading power for false-alarm rates that stay under budget. We also report the boundaries honestly: at $k{=}8$ heavy-tailed corruptions erode the false-alarm budget above $\alpha{=}0.05$, and cross-dataset archives halve the sub-floor power. The resulting story is deliberately focused: CRR is not an MVTec AUROC SOTA or adversarial-robustness method; it is a calibrated, storage-light reliability layer for few-shot inspection systems where raw anomaly scores alone are not sufficient for deployment-quality decisions.
